\subsubsection{Module 5: Quản lý nội dung}

\begin{longtable}{|p{4cm}|p{11cm}|}
\caption{Mô tả Use Case: Xem chi tiết nội dung bị báo cáo}
\label{usecase-5-02} \\
\hline
\endfirsthead
\caption[]{Mô tả Use Case: Xem chi tiết nội dung bị báo cáo} \\
\hline
\endhead
\textbf{Use case ID}       & UC-M5-02 \\ \hline
\textbf{Tên Use case}      & Xem chi tiết nội dung bị báo cáo \\ \hline
\textbf{Các tác nhân}      & Admin \\ \hline
\textbf{Mô tả}             & Hiển thị toàn bộ ngữ cảnh của nội dung bị user báo cáo để Admin đánh giá tính chính xác của báo cáo. \\ \hline
\textbf{Tiền điều kiện}    & Admin đang xem danh sách báo cáo (UC-M5-01). \\ \hline
\textbf{Hậu điều kiện}     & Báo cáo được giải quyết.\\ \hline
\textbf{Luồng sự kiện chính} & 
\begin{tabular}[t]{@{}p{10.5cm}@{}}
1. Hệ thống tải dữ liệu chi tiết của báo cáo được chọn từ UC-M5-01.\\
2. Hệ thống hiển thị giao diện chi tiết gồm 3 phần:\\
- Nội dung gốc đầy đủ.\\
- Thông tin báo cáo.\\
3. Admin đánh giá nội dung dựa trên quy định cộng đồng.\\
4. Hệ thống cung cấp 2 nhóm hành động:\\
- Nhóm 1: "Không vi phạm / Bỏ qua".\\
- Nhóm 2: "Xử lý vi phạm".\\
5. [Điểm mở rộng] Nếu Admin chọn nhóm "Xử lý vi phạm", hệ thống kích hoạt UC-M5-03.\\
6. Sau khi UC-M5-03 hoàn tất, hệ thống quay lại danh sách báo cáo (UC-M5-01).
\end{tabular} \\ \hline
\textbf{Luồng rẽ nhánh}    & 
\begin{tabular}[t]{@{}p{10.5cm}@{}}
\textbf{A1: Quyết định bỏ qua}\\
1. Tại bước 4, Admin nhận thấy báo cáo sai, chọn nút "Bỏ qua".\\
2. Hệ thống cập nhật trạng thái báo cáo thành "Đã từ chối".\\
3. Hệ thống quay về danh sách.
\end{tabular} \\
\hline
\end{longtable}

\begin{longtable}{|p{4cm}|p{11cm}|}
\caption{Mô tả Use Case: Ghi lại lịch sử vi phạm}
\label{usecase-5-04} \\
\hline
\endfirsthead
\caption[]{Mô tả Use Case: Ghi lại lịch sử vi phạm} \\
\hline
\endhead
\textbf{Use case ID}       & UC-M5-04 \\ \hline
\textbf{Tên Use case}      & Ghi lại lịch sử vi phạm \\ \hline
\textbf{Các tác nhân}      & Không có Primary Actor - Included Use Case \\ \hline
\textbf{Mô tả}             & Sub-use case bắt buộc, được \texttt{<<include>>} bởi các Use Case xử lý nội dung vi phạm. Tự động tạo bản ghi vi phạm liên kết với người dùng để Module 8 sử dụng. \\ \hline
\textbf{Tiền điều kiện}    & 
\begin{tabular}[t]{@{}p{10.5cm}@{}}
1. Use Case cha đã được kích hoạt và xác nhận.\\
2. Hệ thống đã xác định được ID người dùng vi phạm.    
\end{tabular} \\ \hline
\textbf{Hậu điều kiện}     & 
\begin{tabular}[t]{@{}p{10.5cm}@{}}
1. (Thành công): Bản ghi \texttt{VIOLATION\_RECORDS} mới được tạo trong CSDL.\\
2. (Thất bại): Lỗi được ghi nhận, Use Case cha bị dừng và thông báo lỗi.
\end{tabular} \\ \hline
\textbf{Luồng sự kiện chính} & 
\begin{tabular}[t]{@{}p{10.5cm}@{}}
1. Use Case được gọi bởi Use Case cha.\\
2. Hệ thống thu thập thông tin: \texttt{user\_id}, \texttt{content\_id}, \texttt{reason}, \texttt{admin\_id}.\\
3. Hệ thống tạo bản ghi mới trong bảng \texttt{VIOLATION\_RECORDS}.\\
4. Use Case kết thúc thành công, trả quyền kiểm soát về Use Case cha.    
\end{tabular} \\ \hline
\textbf{Luồng ngoại lệ} &
\begin{tabular}[t]{@{}p{10.5cm}@{}}
\textbf{E1: Không tìm thấy User ID}\\
1a. Tại bước 2, hệ thống không xác định được user\_id.\\
1b. Use Case thất bại. Use Case cha bị dừng và hiển thị lỗi: "Không thể xử lý vi phạm, không tìm thấy tác giả."\\
\textbf{E2: Lỗi hệ thống}\\
2a. Tại bước 3, hệ thống gặp lỗi khi tạo bản ghi.\\
2b. Use Case thất bại. Use Case cha bị dừng và hiển thị lỗi: "Ghi nhận lịch sử vi phạm thất bại. Đã xảy ra lỗi hệ thống."
\end{tabular} \\ 
\hline
\end{longtable}

\begin{longtable}{|p{4cm}|p{11cm}|}
\caption{Mô tả Use Case: Xem chi tiết nội dung bị gắn cờ}
\label{usecase-5-06} \\
\hline
\endfirsthead
\caption[]{Mô tả Use Case: Xem chi tiết nội dung bị gắn cờ} \\
\hline
\endhead
\textbf{Use case ID}       & UC-M5-06 \\ \hline
\textbf{Tên Use case}      & Xem chi tiết nội dung bị gắn cờ \\ \hline
\textbf{Các tác nhân}      & Admin\\ \hline
\textbf{Mô tả}             & Hiển thị nội dung bị hệ thống tự động nghi ngờ vi phạm. Hệ thống cần làm nổi bật từ khóa hoặc yếu tố gây vi phạm. \\ \hline
\textbf{Tiền điều kiện}    & Admin đang xem danh sách gắn cờ tự động (UC-M5-05). \\ \hline
\textbf{Hậu điều kiện}     & Nội dung được duyệt hoặc chuyển sang xử lý phạt. \\ \hline
\textbf{Luồng sự kiện chính} & 
\begin{tabular}[t]{@{}p{10.5cm}@{}}
1. Admin chọn một nội dung từ danh sách gắn cờ.\\
2. Hệ thống hiển thị chi tiết:\\
- Nội dung gốc với từ khóa vi phạm được làm nổi bật.\\
- Lý do hệ thống gắn cờ.\\
3. Admin đánh giá.\\
4. Hệ thống hiển thị tùy chọn: "Không vi phạm" \space hoặc "Xác nhận vi phạm".    
\end{tabular} \\ \hline
\textbf{Luồng rẽ nhánh} &
\begin{tabular}[t]{@{}p{10.5cm}@{}}
\textbf{A1: Xác nhận vi phạm}\\
1. Admin nhấn "Xử lý" \space $\rightarrow$ Kích hoạt UC-M5-03.\\
\textbf{A2: Không vi phạm}\\
1. Admin nhấn "Cho phép hiển thị" \space $\rightarrow$ Hệ thống gỡ cờ, nội dung được public, xóa khỏi hàng đợi.
\end{tabular} \\
\hline
\end{longtable}
